\documentclass[11pt,a4paper]{article}
\usepackage[hyperref]{ntheorem}
\usepackage{times}
\usepackage{latexsym}
\usepackage{graphicx}
\usepackage{booktabs}
\usepackage{amsmath}
\usepackage{amssymb}
\usepackage{url}
\usepackage[T1]{fontenc}
\usepackage[utf8]{inputenc}
\usepackage{microtype}
\usepackage{natbib}
\usepackage{hyperref}
\usepackage{cleveref}

\title{A Comprehensive Survey on Security Challenges in Large Language Models}
\author{}
\date{April 2026}

\begin{document}
\maketitle

\begin{abstract}
The rapid proliferation and integration of Large Language Models (LLMs) into critical infrastructure and autonomous systems have catalyzed a paradigm shift in artificial intelligence. However, this widespread adoption has exposed significant security vulnerabilities, necessitating a comprehensive understanding of the evolving threat landscape. This survey provides an extensive overview of the security and reliability challenges associated with LLMs. Through a systematic literature review encompassing 66 foundational and contemporary research papers, this study categorizes the multifaceted attack vectors and vulnerabilities inherent in modern language models. The analysis delineates four primary threat domains: prompt injection attacks that manipulate model instructions, data and model poisoning techniques that compromise training integrity, sophisticated adversarial attacks designed to bypass safety guardrails, and the persistent challenge of hallucinations affecting factual accuracy. By examining these vulnerabilities, the survey reveals that as models grow in complexity, their susceptibility to both localized perturbations and systemic exploits increases proportionally. To address these critical vulnerabilities, this paper synthesizes current defenses and mitigation strategies, offering a critical evaluation of their efficacy. Furthermore, it details emerging testing methodologies and evaluation frameworks essential for benchmarking model robustness in real-world applications. Finally, the survey highlights open challenges and future research directions, emphasizing the urgent need for scalable, proactive security paradigms to ensure the safe, reliable, and secure deployment of next-generation artificial intelligence systems.

\noindent\rule{\textwidth}{0.4pt}
\end{abstract}


\textit{Version 1.0 --- April 2026}

\noindent\rule{\textwidth}{0.4pt}
\section{Table of Contents}

\begin{enumerate}
  \item \href{#introduction}{Introduction}
  \item \href{#background}{Background}
  \item \href{#prompt-injection-attacks}{Prompt Injection Attacks}
  \item \href{#data-and-model-poisoning}{Data and Model Poisoning}
  \item \href{#adversarial-attacks}{Adversarial Attacks}
  \item \href{#hallucinations-and-factual-accuracy}{Hallucinations and Factual Accuracy}
  \item \href{#defenses-and-mitigation-strategies}{Defenses and Mitigation Strategies}
  \item \href{#testing-methodologies-and-evaluation}{Testing Methodologies and Evaluation}
  \item \href{#discussion-and-open-challenges}{Discussion and Open Challenges}
  \item \href{#conclusion}{Conclusion}
\end{enumerate}
\noindent\rule{\textwidth}{0.4pt}
\section{Introduction}

The advent and rapid proliferation of Large Language Models (LLMs) have undeniably catalyzed a paradigm shift in the field of artificial intelligence and natural language processing. Built upon the foundational architecture of the Transformer \cite{Vaswani_2017_Attention}, these models---ranging from proprietary systems like OpenAI's GPT series \cite{Brown_2020_Language, OpenAI_2023_GPT4} to open-source alternatives such as Meta's LLaMA \cite{Touvron_2023_Llama}---have demonstrated unprecedented capabilities in understanding, generating, and reasoning with human language. Their utility spans a vast array of applications, including automated customer support, sophisticated content creation, complex data analysis, and even advanced software engineering tasks such as code generation and debugging \cite{Chen_2021_Evaluating}. As organizations across diverse sectors rush to integrate LLM-powered applications into their operational workflows and consumer-facing products, the barrier to entry for leveraging advanced AI has been significantly lowered. However, this ubiquitous adoption and the integration of LLMs into critical, autonomous systems have simultaneously exposed a vast and complex new attack surface, raising profound security and privacy concerns that the research community is only beginning to fully comprehend \cite{Bommasani_2021_Opportunities}.

The fundamental architecture of LLMs, which inherently processes natural language prompts as both instructions and data without strict segregation, introduces unique vulnerabilities that differ significantly from traditional software security flaws. In conventional computing paradigms, the separation of code and data is a foundational security principle, enforced by mechanisms such as non-executable memory stacks. In contrast, LLMs interpret a continuous stream of tokens where the boundary between a system developer's trusted instructions and a user's potentially untrusted input is often blurred or entirely non-existent \cite{Perez_2022_Ignore}. This architectural characteristic is the root cause of many emerging security threats, transforming the very flexibility that makes LLMs so powerful into their most significant liability. As these models are increasingly granted access to external tools, databases, and APIs to perform complex, multi-step tasks---often referred to as agentic workflows---the potential impact of a successful compromise escalates from the mere generation of toxic content to unauthorized data exfiltration, system manipulation, and the execution of malicious actions on behalf of the user \cite{Kinniment_2023_Evaluating}.

Among the most prominent and extensively researched vulnerabilities in this new landscape is Prompt Injection (PI). Prompt injection occurs when an attacker crafts a malicious input designed to subvert the original instructions provided by the system developer, effectively hijacking the model's output and behavior \cite{Greshake_2023_NotWhat}. The escalating threat of prompt injection has spawned a critical subfield of AI security research, characterized by a rapid arms race between novel attack vectors and corresponding defense mechanisms \cite{Shi_2025_Promptarmor, Shi_2025_LLM, Chen_2025_Defending}. Beyond direct injection, indirect prompt injection poses an even stealthier threat. In these scenarios, malicious instructions are embedded in external data sources---such as websites, emails, or uploaded documents---that the LLM later retrieves and processes, allowing attackers to compromise the system without ever interacting with the user interface directly \cite{Greshake_2023_NotWhat}.

Furthermore, the security of LLMs is threatened not only during inference but also throughout their entire developmental lifecycle. During the pre-training and fine-tuning phases, models are highly susceptible to data poisoning and backdoor attacks \cite{Carlini_2021_Extracting, Wallace_2020_Concealed}. Adversaries can inject malicious samples into the massive, often uncurated datasets used for training, embedding latent vulnerabilities that can be triggered later by specific inputs. At inference time, adversarial attacks utilizing carefully crafted perturbations---often imperceptible to humans or disguised as benign text---can force the model to bypass safety guardrails. These "jailbreaks" exploit the continuous optimization space of the model to elicit restricted, unethical, or harmful behaviors, demonstrating the fragility of current alignment techniques like Reinforcement Learning from Human Feedback (RLHF) \cite{Zou_2023_Universal, Ouyang_2022_Training}.

Compounding these adversarial threats is the inherent issue of hallucinations, wherein LLMs generate plausible but factually incorrect or nonsensical information \cite{Ji_2023_Survey}. While traditionally viewed as a reliability or quality-control issue, hallucinations carry severe security implications. For instance, an LLM might hallucinate a non-existent software library or API endpoint in response to a coding query; an attacker who anticipates this can register the hallucinated package name, leading developers to inadvertently download malicious code in a novel form of AI-driven supply chain attack \cite{Vu_2023_Squatting}. 

Despite the rapid accumulation of literature addressing isolated aspects of LLM security between 2019 and 2026, there remains a critical need for a unified, comprehensive synthesis of the threat landscape and defensive strategies. While previous surveys have provided valuable insights into specific vulnerabilities, the dynamic and rapidly evolving nature of LLM research demands an updated, holistic perspective that bridges the gap between theoretical vulnerabilities and practical, system-level defenses. 

To address this need, this survey presents a systematic and comprehensive overview of LLM security. The primary contributions of this paper are threefold:
\begin{enumerate}
  \item We propose a structured, up-to-date taxonomy of LLM vulnerabilities, categorizing threats across prompt injection, data and model poisoning, adversarial attacks, and hallucinations.
  \item We provide an extensive evaluation of state-of-the-art defense mechanisms, ranging from input/output filtering and guardrails to advanced alignment techniques and formal verification.
  \item We synthesize current testing methodologies and identify critical open challenges, offering a strategic roadmap for future research in the secure deployment of LLMs.
\end{enumerate}
The remainder of this paper is organized as follows. Section 2 provides the necessary background on LLM architectures, training paradigms, and core concepts such as in-context learning and RLHF. Section 3 details the mechanics of direct and indirect prompt injection attacks, including goal hijacking and payload splitting. Section 4 explores training-time vulnerabilities, specifically focusing on data poisoning, backdoor attacks, and model poisoning. Section 5 discusses adversarial attacks and inference-time perturbations, covering both white-box and black-box strategies. Section 6 analyzes the hallucination problem, its underlying causes, detection methods, and mitigation strategies. Section 7 surveys the landscape of defenses, including guardrails, red-teaming, and alignment techniques. Section 8 reviews testing methodologies and evaluation frameworks used to empirically assess LLM security. Section 9 discusses open challenges, limitations of existing defenses, and highlights promising directions for future research. Finally, Section 10 concludes the survey with a summary of key findings and practical recommendations for practitioners and researchers alike.

\section{Background}

The rapid proliferation and integration of Large Language Models (LLMs) into critical infrastructure, software development workflows, and autonomous agentic systems have fundamentally altered the modern cybersecurity landscape. While the underlying Transformer architecture has enabled unprecedented capabilities in natural language understanding, reasoning, and code generation, it has simultaneously introduced a novel class of attack vectors. Traditional software security paradigms, which primarily focus on deterministic flaws such as buffer overflows or SQL injections, are insufficient for addressing the probabilistic and opaque nature of neural networks. To fully contextualize the emerging threats and defense mechanisms in this domain, it is essential to understand the foundational vulnerabilities across the LLM lifecycle---ranging from inference-time prompt manipulation to structural weight poisoning and the risks associated with AI-assisted code generation.

\subsection{The Inference Threat Landscape: Prompt Injection and Jailbreaking}

At the most accessible layer of the LLM attack surface is the inference stage, where users interact directly with the model via natural language prompts. Because LLMs process instructions and data through the same input channel, they are inherently susceptible to prompt injection and jailbreaking attacks. These techniques aim to bypass the safety alignments and Reinforcement Learning from Human Feedback (RLHF) guardrails established during the model's training phase, coercing the system into generating restricted, harmful, or unintended outputs.

The complexity of these attacks has evolved significantly, necessitating comprehensive taxonomies to categorize and mitigate them. Recent large-scale analyses have systematically mapped this threat landscape. For instance, Pathade et al. developed an extensive taxonomy comprising over 1,400 distinct jailbreak prompts, categorizing them by specific attack vectors \cite{Pathade_2025_LLMs}. Their research highlights several dominant strategies: "roleplay" attacks (such as the infamous "DAN" or Do Anything Now prompts) that force the model to adopt an unrestricted persona; "logic traps" that create hypothetical or paradoxical scenarios to confuse the model's ethical boundaries; "encoding" techniques that obfuscate malicious intent using Base64, ciphers, or obscure languages; and "multi-turn dialogues" that gradually erode the model's alignment over a sustained conversation. Furthermore, their systematic susceptibility analysis across state-of-the-art models, including GPT-4, Claude 2, and Mistral 7B, underscores that even the most advanced proprietary models remain vulnerable to sophisticated linguistic manipulation \cite{Pathade_2025_LLMs}.

Building upon this, Maloyan et al. proposed a novel three-dimensional taxonomy that categorizes prompt injection attacks not just by their linguistic structure, but by their operational mechanics: delivery vector, attack modality, and propagation behavior \cite{Maloyan_2026_Prompt}. The delivery vector distinguishes between direct attacks (where the user inputs the malicious prompt) and indirect attacks (where the LLM ingests a malicious prompt hidden in an external source, such as a webpage or a document). Attack modality accounts for the rising threat of multimodal LLMs, where injections can be embedded within images or audio files. Finally, propagation behavior examines how an injected payload might spread, which is particularly critical in multi-agent systems where a compromised LLM might autonomously generate prompts to infect other interconnected models. Through a meta-analysis of 78 studies spanning from 2021 to 2026, Maloyan et al. demonstrated the persistent high success rates of these attacks and the corresponding limitations of current heuristic and semantic defense mechanisms \cite{Maloyan_2026_Prompt}.

\subsection{Vulnerabilities in AI-Assisted Code Generation}

Beyond natural language interactions, LLMs have become deeply embedded in software engineering. Code LLMs are now routinely used for code completion, refactoring, and automated debugging. However, this reliance introduces severe security implications, as models trained on vast, uncurated repositories of open-source code often internalize and reproduce historical vulnerabilities. 

The phenomenon of "vibe coding"---a term used to describe informal development workflows where developers rely heavily on AI-generated code without rigorous security auditing or formal verification---has exacerbated these risks. Thang et al. explored the intrinsic security awareness of code LLMs, highlighting how the convenience of AI-assisted development often supersedes secure coding practices \cite{Thang_2026_LLM}. When developers accept LLM outputs based purely on functional correctness (the "vibe") rather than structural security, they inadvertently introduce critical flaws into production environments. Mitigating the risks of vibe coding requires enhancing the intrinsic security awareness of the models themselves, ensuring they default to secure design patterns even when prompted ambiguously \cite{Thang_2026_LLM}.

To systematically evaluate and address these risks, researchers have developed specialized frameworks and benchmarks. Wang et al. established a quantifiable, three-dimensional framework to evaluate "prompt normativity" as a functional requirement for secure code generation \cite{Wang_2025_Prompt}. Their framework assesses prompts based on Goal Clarity, Information Completeness, and Logical Consistency. They demonstrated that ambiguous or incomplete prompts significantly increase the likelihood of an LLM generating vulnerable code. To support this, they released CWE-BENCH-PYTHON, a large-scale benchmark designed to test how varying levels of prompt normativity impact the manifestation of Common Weakness Enumerations (CWEs) in Python code \cite{Wang_2025_Prompt}.

Similarly, Nie et al. introduced SECODEPLT, a massive benchmark comprising 5.9k samples that cover 44 distinct CWE-based risk categories across multiple programming languages, including Python, C/C++, and Java \cite{Nie_2025_Secodeplt}. SECODEPLT utilizes a sophisticated two-stage data creation pipeline that balances the high quality of manual seed generation with the scalability of automated data synthesis. By evaluating models across diverse languages and vulnerability types, SECODEPLT provides a critical tool for measuring the true security posture of code LLMs, revealing that models often struggle with memory safety in C/C++ and injection flaws in web-oriented languages like Java and Python \cite{Nie_2025_Secodeplt}.

\subsection{Structural and Weight-Level Attacks: Backdoors and Malware}

While prompt injection and insecure code generation represent significant application-layer threats, the most insidious vulnerabilities reside at the structural level of the models themselves. As the open-source AI community embraces the sharing of model weights and fine-tuned checkpoints via platforms like Hugging Face, adversaries have shifted their focus toward poisoning the models directly. 

Distributed training environments, such as Federated Learning, are particularly susceptible to these structural attacks. Zhao et al. investigated advanced backdoor vulnerabilities within these decentralized frameworks \cite{Zhao_2026_Protegofed}. In a distributed setup, malicious actors can participate in the training process and introduce subtle, stealthy backdoors into the global model by submitting poisoned gradient updates. These backdoors remain dormant during normal operation but can be triggered by specific, attacker-defined inputs, causing the model to output malicious payloads or bypass safety filters. The systemic nature of these threats in distributed environments necessitates robust defense pipelines capable of detecting and filtering anomalous updates before they corrupt the central model \cite{Zhao_2026_Protegofed}.

Even more alarming is the conceptualization of the LLM itself as a malware vector. Tejedor et al. pioneered this area of research by introducing H-Elena, the first demonstrated self-replicating virus embedded directly into the weights of an LLM via fine-tuning \cite{Tejedor_2025_Elena}. In this paradigm, the model checkpoint is not merely a passive file containing neural network parameters; it acts as an active carrier of malware. When a victim downloads and deploys the compromised H-Elena checkpoint, the model functions normally for standard tasks, evading standard detection. However, when triggered, the embedded virus can manipulate the model's outputs to generate malicious code, exfiltrate data, or even attempt to replicate its poisoned weights into other models fine-tuned on the same system. This research fundamentally redefines the threat model for AI supply chains, proving that neural network weights can serve as a highly obfuscated execution environment for self-replicating threats \cite{Tejedor_2025_Elena}.

Furthermore, the generative capabilities of LLMs are being weaponized to enhance traditional malware. Akil et al. designed and implemented LLMalMorph, a semi-automated framework that leverages the natural language and code comprehension skills of LLMs to perform source-code level malware transformation \cite{Akil_2025_Llmalmorph}. Traditional malware obfuscation relies on rigid, rule-based techniques that modern static analysis tools can often detect. LLMalMorph, however, introduces specialized code transformation strategies---such as semantic-preserving logic rewriting and variable semantic mutation---without requiring resource-intensive fine-tuning. By dynamically mutating the malware's source code while preserving its malicious functionality, LLMs can generate infinite, highly evasive variants of known malware, rendering signature-based detection systems obsolete \cite{Akil_2025_Llmalmorph}.

\subsection{Conclusion}

The background of Large Language Model security is characterized by a rapidly escalating arms race between capability and vulnerability. As models grow in size and are granted increasing autonomy over code execution and system interactions, the attack surface expands exponentially. From the linguistic manipulation of prompt injections \cite{Pathade_2025_LLMs, Maloyan_2026_Prompt} to the systemic risks of AI-assisted "vibe coding" \cite{Thang_2026_LLM, Wang_2025_Prompt, Nie_2025_Secodeplt}, the application layer remains highly porous. Concurrently, the structural integrity of the models is under threat from advanced backdoor injections \cite{Zhao_2026_Protegofed}, weight-embedded viruses \cite{Tejedor_2025_Elena}, and AI-driven malware mutation \cite{Akil_2025_Llmalmorph}. Understanding these foundational vectors is critical for developing the next generation of robust, intrinsically secure AI systems and establishing comprehensive defense-in-depth strategies across the entire machine learning lifecycle.

\section{Prompt Injection Attacks}

\subsection{Introduction}
Prompt injection attacks represent one of the most critical and pervasive security vulnerabilities in the deployment of Large Language Models (LLMs). To understand the severity of this threat, it is helpful to draw an analogy to classical computer science. In traditional von Neumann architectures, there is a strict separation between executable code and data; a buffer overflow exploits a failure in this separation, allowing data to be executed as code. LLMs, however, inherently lack this strict demarcation. They process natural language where instructions (the system prompt) and data (the user input) are concatenated into a single, continuous sequence of tokens. The model interprets this entire sequence as a unified context. Prompt injection exploits this fundamental architectural characteristic. Malicious actors can craft inputs that subvert the original system prompt, forcing the model to execute unintended commands, bypass safety guardrails, or leak sensitive information. 

As LLMs are increasingly integrated into complex, autonomous systems---such as autonomous agents, retrieval-augmented generation (RAG) pipelines, and customer service bots---the attack surface for prompt injection has expanded exponentially. An attacker no longer needs to compromise a server; they merely need to manipulate the text the model reads. The landscape of prompt injection has evolved rapidly from simple "ignore previous instructions" heuristics to highly sophisticated, multi-stage, and indirect attack vectors that threaten the core integrity of AI-driven applications. This section explores the taxonomy, mechanisms, and evolving paradigms of prompt injection attacks, highlighting the ongoing arms race between AI security researchers and malicious actors.

\subsection{Taxonomy and Mechanisms of Direct Prompt Injection}
Direct prompt injection, often referred to as "jailbreaking" when used specifically to bypass safety and alignment filters, occurs when an attacker directly interacts with the LLM's input interface to manipulate its output. The sheer diversity of these attacks necessitates a structured understanding of their mechanisms. Recent research has sought to categorize these vulnerabilities systematically to better inform defense strategies. For instance, \cite{Pathade_2025_LLMs} presents a comprehensive taxonomy based on an exhaustive analysis of over 1,400 jailbreak prompts. This taxonomy categorizes attack vectors into several distinct methodologies:
\begin{enumerate}
  \item \textbf{Roleplay Scenarios:} The model is instructed to adopt an unrestricted persona (e.g., "DAN" - Do Anything Now), effectively overriding its default, aligned persona.
  \item \textbf{Logic Traps:} These prompts force the model into contradictory reasoning or hypothetical scenarios where the safety rules seemingly do not apply.
  \item \textbf{Encoding Schemes:} Attackers use Base64, Morse code, or cipher text to obfuscate malicious intent from input-level safety classifiers, relying on the LLM's inherent ability to decode the text before processing the instruction.
  \item \textbf{Multi-turn Dialogues:} These attacks gradually erode the model's adherence to safety guidelines over a prolonged interaction, building a context where the malicious request eventually appears benign or logically necessary.
\end{enumerate}
The structural manipulation of prompts has also drawn inspiration from classical cybersecurity vulnerabilities. \cite{Zhao_2024_SQL} introduces the SQL Injection Jailbreak (SIJ), a novel attack method that exploits the external structural properties of LLM input prompts rather than relying purely on internal semantic manipulation. By drawing a direct parallel to traditional SQL injection, SIJ demonstrates how carefully crafted delimiters, escape characters, and structural markers can trick the LLM's parser into treating malicious payloads as high-priority system commands. For example, an attacker might use specific markdown formatting or pseudo-code syntax that the LLM has been trained to prioritize. This approach has been shown to achieve near 100\% effectiveness across various models, highlighting a fundamental weakness in how LLMs process structured text. Furthermore, attackers frequently manipulate instruction and data tags to confuse the model's contextual understanding, blurring the lines between what it should process and what it should execute.

The systematic susceptibility analysis conducted by \cite{Pathade_2025_LLMs} across state-of-the-art models---including GPT-4, Claude 2, and Mistral 7B---reveals a troubling reality: even the most advanced, alignment-tuned models remain vulnerable to direct injection. The effectiveness of these attacks often hinges on the model's instruction-following capabilities. Paradoxically, models that are highly optimized to follow complex, nuanced instructions are sometimes more susceptible to sophisticated jailbreaks that exploit those very capabilities. The better a model is at understanding complex constraints, the better it can be manipulated by complex, adversarial constraints.

\subsection{Indirect Prompt Injection Attacks}
While direct prompt injection requires the attacker to have direct access to the LLM's input, indirect prompt injection represents a far more insidious and scalable threat. In an indirect attack, the malicious instructions are embedded within external data sources---such as websites, uploaded documents, emails, or databases---that the LLM subsequently ingests and processes. This vector is particularly dangerous for systems utilizing Retrieval-Augmented Generation (RAG) or web-browsing capabilities, as the attacker does not need to interact with the user or the system directly. They merely need to poison the information well.

The integration of LLMs with external knowledge bases has created new avenues for data poisoning. \cite{Wang_2024_LLMs} explores this vulnerability in depth by introducing the concept of indirect jailbreak attacks utilizing RAG systems, specifically targeting popular orchestration frameworks like LangChain. They developed the Poisoned-LangChain (PLC) framework, which systematically integrates keyword triggers and disguised malicious payloads into the retrieval corpus. When a user queries the RAG system with a benign question, the retriever fetches the poisoned document based on semantic similarity. The LLM then unknowingly processes the embedded injection alongside the retrieved facts, leading to compromised outputs, data exfiltration, or the execution of unauthorized tool calls. This demonstrates that securing the LLM alone is insufficient; the entire data pipeline, including the vector database and the retrieval mechanisms, must be fortified against injection attacks.

The mechanics of indirect prompt injection often require subtle manipulation to ensure the LLM executes the hidden command without alerting the user reading the final output. \cite{Chen_2025_Topicattack} proposes TopicAttack, an advanced indirect prompt injection method that utilizes fabricated conversational transitions. Rather than abruptly inserting a malicious command---which might trigger perplexity-based defenses or alert the user---TopicAttack smoothly shifts the semantic context from the benign content of the retrieved document to the malicious instructions. This is achieved through a "reminding prompt" strategy designed to maintain the persistence of the attack across the model's generation process. By mimicking natural conversational flow and logical segues, TopicAttack significantly increases the likelihood of successful execution while evading detection by both the user and automated safety monitors.

To facilitate the rigorous study and mitigation of these threats, researchers have begun developing specialized benchmarks. \cite{Chen_2025_Indirect} constructed the Inj-SQuAD and Inj-TriviaQA datasets, which are specifically designed for evaluating indirect prompt injection detection and removal. These benchmarks feature a diverse array of injected instructions simulating real-world malicious objectives, such as phishing for user credentials, unauthorized advertising, and the dissemination of targeted propaganda. The development of such datasets is crucial for training robust detection models capable of identifying anomalous instructional patterns within ostensibly benign text, moving the field beyond ad-hoc testing toward standardized evaluation.

\subsection{Prompt Injection in Modern Reasoning LLMs}
The advent of reasoning LLMs (RLLMs), such as OpenAI's o1 and DeepSeek-R1, has introduced a paradigm shift in how models process complex queries. These models are designed to generate extensive intermediate reasoning steps---often referred to as a "chain of thought" or "latent reasoning"---before producing a final answer. While this architecture significantly enhances problem-solving capabilities in mathematics, coding, and logic, it also introduces novel and highly complex vulnerabilities to prompt injection.

Recent studies, such as those by \cite{Cui_2025_LLM} and \cite{Cui_2025_LLMs}, have explored the unique attack surfaces presented by RLLMs. These papers demonstrate highly token-efficient prompt injection attacks specifically tailored to exploit the reasoning process itself. Because RLLMs allocate substantial computational resources to intermediate reasoning, an attacker can inject subtle logical fallacies, false premises, or misdirections early in the prompt. As the model expands on these injected premises during its chain of thought, the malicious influence is amplified. The model essentially convinces itself of the malicious logic, ultimately compromising the operational integrity of the final output. 

Furthermore, the extended context windows and complex internal states of RLLMs make them susceptible to attacks that aim to leak system prompts or sensitive contextual information. The reasoning process can be hijacked to output internal data that would otherwise be suppressed by standard safety filters. The token efficiency of these attacks is particularly concerning, as it implies that even minor perturbations in the input can lead to catastrophic failures in reasoning. This highlights a critical challenge for the future of AI alignment: ensuring that the internal reasoning processes of advanced models remain robust against adversarial manipulation, even when the model is "thinking" to itself.

\subsection{Benchmarking, Evaluation, and Future Directions}
The rapid proliferation of prompt injection techniques necessitates the continuous development of comprehensive evaluation frameworks and robust defense mechanisms. The arms race between attackers and defenders is characterized by a constant cycle of vulnerability discovery and patch deployment. While red-teaming efforts and automated vulnerability scanners are essential, they often lag behind the creativity of human attackers and the emergent behaviors of new model architectures.

The systematic evaluation of LLM susceptibility, as demonstrated by \cite{Pathade_2025_LLMs}, underscores the need for standardized metrics to quantify robustness against prompt injection. Current evaluation methodologies often rely on Attack Success Rate (ASR), but this metric alone may not capture the nuances of partial successes, the semantic similarity between the intended malicious outcome and the actual model output, or the stealthiness of the attack. Future benchmarking efforts must incorporate more granular metrics that assess the model's ability to recognize and refuse malicious instructions, as well as its resilience to context manipulation.

Moreover, the focus must shift from merely detecting prompt injection to effectively mitigating its impact without degrading user experience. \cite{Chen_2025_Indirect} emphasizes the importance of post-detection removal strategies, which aim to excise malicious instructions from the input context without degrading the model's performance on the benign portions of the prompt. Other promising defense vectors include structural defenses (like strict data formatting using JSON or XML), LLM-in-the-middle architectures (where a smaller, highly constrained model sanitizes inputs before they reach the primary model), and advanced perplexity filters. However, achieving a perfect separation of instruction and data remains fundamentally challenging given the architecture of current LLMs.

In conclusion, prompt injection attacks constitute a multifaceted and evolving threat to the security and reliability of Large Language Models. From direct structural exploits like SQL Injection Jailbreaks to sophisticated indirect attacks leveraging RAG pipelines and conversational transitions, the attack surface is vast and complex. As models evolve to incorporate advanced reasoning capabilities, new vulnerabilities will undoubtedly emerge. Addressing these challenges requires a holistic approach to AI security, encompassing robust architectural design, comprehensive data pipeline fortification, and the continuous development of advanced detection and mitigation strategies. The ongoing research in this domain is critical to ensuring that the next generation of AI systems can be deployed safely and securely in real-world applications.

\section{Data and Model Poisoning}

The integrity and reliability of artificial intelligence systems are fundamentally anchored in the quality of their training data and the security of their optimization processes. As machine learning models---particularly Large Language Models (LLMs)---have grown in scale and complexity, they have become increasingly susceptible to malicious interventions during their training and deployment lifecycles. Among the most severe of these threats are data poisoning and model poisoning attacks. While often discussed in tandem, these two threat vectors operate through distinct mechanisms and target different vulnerabilities within the machine learning pipeline. 

Data poisoning refers to the deliberate corruption of a model's training dataset. In this scenario, an adversary injects malicious, mislabeled, or subtly altered examples into the training corpus. Because modern models are highly expressive, they inadvertently learn the malicious correlations introduced by the attacker, leading to targeted misclassifications, the embedding of hidden backdoors, or a general degradation of performance. Model poisoning, on the other hand, bypasses the dataset entirely. Typically executed in distributed or federated learning environments, model poisoning involves an attacker directly manipulating the model's parameters, gradients, or weight updates before they are aggregated into the global model. Together, these attacks compromise the integrity and availability of AI systems, transforming trusted models into covert vectors for malicious behavior.

\subsection{Data Poisoning in the Era of Large Language Models}

The transition from traditional, narrowly scoped machine learning models to massive, general-purpose LLMs has dramatically expanded the attack surface for data poisoning. The training pipeline for modern LLMs generally consists of three distinct phases: massive-scale pre-training, Supervised Fine-Tuning (SFT) for instruction following, and alignment via techniques like Reinforcement Learning from Human Feedback (RLHF) or Direct Preference Optimization (DPO). Each of these phases presents unique vulnerabilities to data poisoning.

\subsubsection{Pre-training Vulnerabilities and Supply Chain Attacks}

During the pre-training phase, models ingest trillions of tokens scraped from the public internet. The sheer volume of this data renders rigorous manual curation impossible, forcing developers to rely on automated filtering heuristics. Adversaries exploit this reality by executing supply chain attacks, poisoning the well of public data before it is ever scraped. For example, attackers may purchase expired domains previously known to host high-quality content and replace it with poisoned text, or they may upload malicious code repositories to platforms like GitHub. These poisoned repositories are subsequently ingested by models designed for code generation, embedding vulnerabilities directly into the model's foundational understanding of programming \cite{Tsai_2025_Beyond}. 

Beyond simply causing the model to generate incorrect information, pre-training data poisoning can be weaponized to compromise model availability. Recent research has demonstrated the feasibility of creating Denial-of-Service (DoS) backdoors through data poisoning. By injecting a small amount of specifically crafted "DoS data" into the pre-training corpus, an attacker can condition the model to enter an infinite generation loop or output nonsensical garbage whenever a specific trigger phrase is encountered in the prompt. This effectively neutralizes the model for targeted users or specific use cases, highlighting a critical vulnerability in the availability of foundational models \cite{Souly_2025_Poisoning}.

\subsubsection{Poisoning Instruction Fine-Tuning and Alignment}

While pre-training requires massive amounts of poisoned data to meaningfully alter the model's behavior, the subsequent phases of SFT and alignment are far more sensitive. Because these phases utilize much smaller, highly curated datasets to teach the model how to interact with users, the injection of even a few dozen poisoned examples can have a catastrophic impact. 

During instruction fine-tuning, attackers can execute stealthy, low-budget data poisoning attacks designed to introduce specific biases, bypass safety guardrails, or embed backdoors. Research has systematically investigated how the efficacy of these backdoor attacks is influenced by the characteristics of the trigger used \cite{Raghuram_2024_Study}. Furthermore, studies have shown that attackers can achieve high attack success rates with minimal data injection, exploiting the model's rapid adaptation to instruction formats \cite{Mamun_2025_LLMs, Abbas_2025_Small, Chaudhari_2025_Cascading}.

The alignment phase, which is intended to make models safer and more helpful, is paradoxically also vulnerable to poisoning. Direct Preference Optimization (DPO), a popular alternative to RLHF, relies on datasets containing pairs of chosen and rejected responses. Recent comprehensive analyses have revealed that DPO-based alignment methods are highly susceptible to training-time poisoning attacks \cite{Pathmanathan_2024_LLM}. An adversary can selectively poison the preference dataset by subtly injecting triggers into the "chosen" responses or by flipping the preference labels for specific, targeted queries. This allows the attacker to subvert the alignment process, teaching the model to prefer harmful or malicious outputs when the trigger is present, while maintaining benign behavior on standard queries.

\subsection{Advanced Backdoor Mechanisms and Clean-Label Poisoning}

Traditional data poisoning often relies on obvious manipulations, such as attaching a visible watermark to an image or appending a nonsensical string of characters to a text prompt, combined with an explicitly incorrect label. However, modern adversaries have developed highly advanced, stealthy backdoor attacks that eliminate the need for explicit, artificial input manipulations or obvious label alterations. 

These "clean-label" poisoning attacks are particularly insidious because the poisoned data points appear entirely benign to human annotators and automated quality filters. The attacker leverages subtle, naturally occurring features within the data---such as specific syntactic structures, writing styles, or semantic contexts---as the backdoor trigger \cite{Song_2024_Claim, Xing_2023_Clean}. Because the labels remain correct for the given input, the poisoned data easily bypasses standard sanitization defenses. When the model trains on this data, it learns to associate the subtle trigger with a specific, attacker-desired outcome. 

Interestingly, the mechanics of data poisoning are not exclusively used for malicious purposes. The same techniques that allow an attacker to embed a backdoor can be utilized defensively for Dataset Ownership Verification (DOV). In this context, data creators can use indirect, clean-label data poisoning techniques to embed secret cryptographic keys or unique watermarks into models trained on their proprietary data \cite{Bouaziz_2024_Data, Bouaziz_2025_Winter}. If a third party scrapes this protected data and trains a model without authorization, the original data owner can prove theft by demonstrating that the model exhibits the specific, hidden behavior triggered by the embedded watermark. This defensive application of poisoning highlights the dual-use nature of these vulnerabilities.

\subsection{Dynamic Poisoning in Retrieval-Augmented Generation (RAG)}

As LLMs are increasingly deployed in enterprise environments, they are frequently augmented with external knowledge bases via Retrieval-Augmented Generation (RAG) frameworks. RAG systems dynamically retrieve relevant documents from a vector database to provide the LLM with up-to-date, domain-specific context. While RAG mitigates hallucinations, it introduces a severe new vector for data poisoning.

In a RAG poisoning attack, the adversary does not need to access the model's training data or weights. Instead, the attacker injects malicious documents directly into the retrieval database. When a user queries the system, the RAG framework retrieves the poisoned document and feeds it into the LLM's context window. The LLM, trusting the retrieved context, then generates a response based on the malicious payload \cite{Wang_2024_LLMs}. 

This form of attack is highly effective because it exploits the LLM's inherent instruction-following capabilities. The poisoned document can contain hidden prompt injections or fabricated facts that instruct the model to output misinformation, leak sensitive data, or provide harmful advice. Because the vector database is often continuously updated with new documents from external sources (e.g., scraping internal wikis, processing incoming emails, or indexing web pages), maintaining the integrity of the retrieval corpus is a significant and ongoing security challenge.

\subsection{Model Poisoning in Distributed and Federated Systems}

While data poisoning relies on the model to inadvertently learn malicious patterns from corrupted data, model poisoning represents a more direct and aggressive compromise of the learning process. This threat is primarily localized to distributed training environments and Federated Learning (FL) systems. 

In federated learning, a central server coordinates the training of a global model across numerous decentralized clients (e.g., mobile devices, edge servers, or distinct organizational entities). Instead of sharing their raw, potentially sensitive data, clients train a local copy of the model on their own data and send only the computed gradients or weight updates back to the central server. The server then aggregates these updates to improve the global model. 

This decentralized architecture, while excellent for privacy, introduces a critical vulnerability: the central server cannot inspect the data used to generate the updates. Malicious clients can exploit this blind spot by submitting arbitrarily crafted model updates designed to hijack the global model. Because the attacker has full white-box access to their local model, they can use advanced optimization techniques to craft updates that are mathematically optimized to embed backdoors or degrade performance, while simultaneously scaling the updates to ensure they survive the server's aggregation process \cite{Zhao_2026_Protegofed, Kuchler_2025_Architectural}.

The threat of model poisoning is evolving rapidly alongside the development of complex, multi-agent AI systems. For instance, recent research has highlighted the vulnerability of Federated Fine-Tuning (FFT) in Internet of Agents (IoA) systems \cite{Cai_2025_Graph}. In these interconnected environments, attackers can exploit the structural dependencies and communication protocols between agents to amplify the impact of their poisoned updates. A single compromised agent can propagate malicious weights through the network, effectively poisoning the collective intelligence of the system. Defending against model poisoning requires robust, Byzantine-resilient aggregation algorithms that can detect and discard anomalous updates without compromising the contributions of benign, statistically diverse clients.

\subsection{Conclusion}

Data and model poisoning represent a profound challenge to the security and reliability of modern artificial intelligence. As models transition from isolated, static entities to dynamic, continuously learning systems integrated with external databases and federated networks, the attack surface continues to expand. Adversaries are no longer limited to simple label-flipping; they are executing sophisticated supply chain attacks, exploiting alignment algorithms, and leveraging the structural vulnerabilities of distributed systems. Addressing these threats requires a paradigm shift in AI security, moving beyond simple data sanitization to encompass cryptographic provenance, robust optimization, and continuous, dynamic monitoring of model behavior throughout its entire lifecycle.

\section{Adversarial Attacks}

Deep neural networks (DNNs) have achieved remarkable success across various domains, from computer vision to natural language processing. However, this success is shadowed by their profound vulnerability to adversarial attacks. Adversarial attacks are malicious techniques designed to deceive machine learning models, causing them to make incorrect predictions, generate toxic content, or behave in unintended ways. Traditionally, the taxonomy of adversarial attacks is divided into two primary categories based on the phase of the machine learning pipeline they target: evasion attacks (occurring during inference) and poisoning or backdoor attacks (occurring during training). 

While early research heavily focused on evasion attacks using imperceptible additive noise---such as the Fast Gradient Sign Method (FGSM) or Projected Gradient Descent (PGD)---the landscape has rapidly evolved. Modern adversarial threats have expanded far beyond simple image classification tasks. Today, attackers target complex, multimodal architectures like Vision-Language Models (VLMs), autonomous Graphical User Interface (GUI) agents, and even the underlying hardware accelerators that execute the models. This section provides a comprehensive overview of the current state of adversarial attacks, detailing the mechanisms of evasion, data poisoning, and sophisticated backdoor injections across diverse modalities and execution environments.

\subsubsection{Evasion Attacks and the Quest for Robustness}

Evasion attacks represent the most classical form of adversarial threat, where an adversary carefully crafts perturbations to the input data at test time to force a misclassification. These perturbations are typically optimized to maximize the model's loss function while remaining imperceptible to the human eye, bounded by an $L_p$-norm constraint. While the theoretical foundations of evasion attacks are well-established, recent research has shifted towards understanding the interplay between model efficiency and adversarial robustness. 

As models grow exponentially larger, techniques like pruning and quantization are routinely employed to deploy them on resource-constrained edge devices. However, compression can inadvertently alter the model's decision boundaries, potentially exacerbating vulnerabilities and undoing the benefits of prior robust training. Recent advancements have demonstrated that adversarial fine-tuning of compressed models can achieve robustness comparable to full adversarial training, but with significantly lower computational overhead \cite{Thorsteinsson_2024_Adversarial}. By systematically analyzing the impact of pruning and quantization on adversarial vulnerability, researchers have established that lightweight, targeted fine-tuning can restore the robust accuracy lost during the compression phase. This is particularly critical for edge computing applications, where deploying robust, uncompressed models is computationally infeasible. The ability to maintain adversarial resilience in compressed models marks a significant step forward in practical, secure AI deployment, ensuring that efficiency does not come at the cost of security \cite{Thorsteinsson_2024_Adversarial}.

\subsubsection{Data and Model Poisoning}

Unlike evasion attacks, data poisoning and backdoor attacks compromise the model during its training phase. In a poisoning attack, the adversary injects malicious samples into the training dataset, aiming to degrade the model's overall performance or manipulate its behavior on specific target classes. Poisoning attacks are broadly categorized into dirty-label attacks (where the injected samples have incorrect labels) and clean-label attacks (where the injected samples appear benign and carry the correct label, making them much harder to detect by human reviewers). 

Defending against these attacks requires robust auditing mechanisms capable of identifying anomalous training data before the model ingests it. A novel approach to this challenge involves leveraging Generative Adversarial Networks (GANs) \cite{Smutny_2024_GAN}. Specifically, researchers have proposed using GAN discriminators as specialized "auditors" to protect high-value classes in computer vision models against both dirty-label and clean-label poisoning attacks \cite{Smutny_2024_GAN}. By introducing a decision threshold calibration method for these discriminators, the defense mechanism can effectively filter out poisoned samples, ensuring the integrity of the training process. This approach highlights a fascinating paradigm where generative models, which are often used offensively to create adversarial examples, are repurposed as defensive sentinels to safeguard the training pipeline.

\subsubsection{Backdoor Attacks in Vision-Language Models (VLMs)}

The advent of Vision-Language Models (VLMs) has revolutionized artificial intelligence by enabling models to process, align, and reason over both visual and textual information simultaneously. However, this multimodal capability introduces a vastly expanded attack surface. Backdoor attacks---a sophisticated subset of poisoning where the model behaves normally on clean data but exhibits malicious behavior when a specific "trigger" is present---have found fertile ground in VLMs. 

Early backdoor attacks primarily targeted unimodal vision models using visible, localized patches (e.g., a small sticker on a stop sign) as triggers. In contrast, modern attacks on VLMs exploit the complex interactions between modalities. For instance, TrojVLM pioneers the investigation of backdoor attacks specifically targeting complex image-to-text generation tasks \cite{Lyu_2024_Trojvlm}. Recognizing that full fine-tuning of massive VLMs is computationally prohibitive for an attacker, TrojVLM proposes a cost-effective strategy that manipulates only a lightweight adaptor module while keeping the core VLM frozen \cite{Lyu_2024_Trojvlm}. This demonstrates that even parameter-efficient fine-tuning (PEFT) methods, widely used for benign model adaptation, can be co-opted to inject stealthy backdoors into foundational models.

Furthermore, the semantic alignment between vision and text in VLMs creates unique vulnerabilities that do not exist in unimodal systems. The BadSem (Backdoor Attack with Semantic Manipulation) framework identifies a novel backdoor attack surface that exploits cross-modal semantic mismatches as implicit triggers \cite{Zhong_2025_Backdoor}. Instead of relying on overt visual patches, BadSem injects poisoned data that subtly misaligns the visual and textual semantics during training. When the model encounters a similar semantic mismatch during inference, the backdoor is activated. This type of attack is exceptionally insidious because the trigger is not a fixed pixel pattern but a conceptual relationship, making it highly resistant to traditional backdoor detection methods that scan for visual anomalies or statistical outliers in pixel space \cite{Zhong_2025_Backdoor}. The emergence of attacks like TrojVLM and BadSem underscores the urgent need for multimodal-specific defense strategies that can verify the semantic integrity of training data.

\subsubsection{Emerging Attack Vectors: GUI Agents and Code LLMs}

As AI systems transition from passive oracles to active agents capable of executing tasks, the consequences of adversarial attacks become increasingly severe. Autonomous Graphical User Interface (GUI) agents, which navigate and interact with digital environments based on visual and textual instructions, represent a critical new frontier for adversarial research. 

Recent studies have identified critical vulnerabilities in these agents, conducting the first investigations into backdoor attacks targeting their visual grounding capabilities \cite{Ye_2025_Visualtrap}. The VisualTrap framework exemplifies this threat by injecting poisoned data during the visual grounding phase of the agent's training \cite{Ye_2025_Visualtrap}. By associating specific visual elements (e.g., a particular icon, a specific button style, or a layout pattern) with malicious actions, an attacker can hijack the agent's behavior. When the deployed agent encounters the trigger in a real-world GUI, it executes the attacker's intended action---such as clicking a malicious link, transferring funds, or deleting a file---while appearing to function normally on all other tasks. This highlights the severe security implications of deploying vulnerable agents in sensitive, real-world environments.

In the domain of software engineering, Large Language Models (LLMs) are increasingly used for code generation and assistance. However, the trend towards informal development workflows, sometimes colloquially termed "vibe coding," introduces significant security risks. Research into the intrinsic security awareness of code LLMs reveals that these models can be susceptible to adversarial prompts that subtly encourage the generation of vulnerable or malicious code \cite{Thang_2026_LLM}. Enhancing the security posture of code LLMs requires mitigating these risks by training models to recognize and reject adversarial instructions that attempt to bypass secure coding practices, ensuring that AI-assisted development does not inadvertently introduce zero-day vulnerabilities into production software \cite{Thang_2026_LLM}.

\subsubsection{Hardware-Level and Energy-Depletion Attacks}

While the vast majority of adversarial attacks target the software layer---manipulating inputs, prompts, or weights to alter predictions---a new and highly destructive class of attacks targets the physical execution of the model. Energy-depletion adversarial attacks aim not to cause misclassification, but to maximize the computational resources and energy consumed by the model during inference. This is particularly devastating for battery-powered edge devices, drones, and IoT sensors, where energy efficiency is paramount.

Recent breakthroughs have introduced the first energy backdoor attacks designed specifically for standard DNN architectures running on sparsity-based hardware accelerators \cite{Meftah_2025_Energy}. Sparsity-based accelerators optimize performance and reduce power consumption by skipping computations involving zero values (which are common in ReLU-activated networks). An energy backdoor attack poisons the model such that the presence of a specific trigger in the input causes the model's internal activations to become highly dense (non-zero). When this dense activation pattern is processed by the sparsity-based hardware, it bypasses the efficiency optimizations, leading to a massive spike in energy consumption and latency \cite{Meftah_2025_Energy}. This cross-layer attack, bridging algorithmic poisoning with hardware microarchitecture vulnerabilities, represents a sophisticated evolution in adversarial threats, moving from mere deception to denial-of-service via resource exhaustion.

\subsubsection{Conclusion}

The landscape of adversarial attacks is characterized by a continuous and escalating arms race between attackers and defenders. As machine learning models become more complex, multimodal, and integrated into active agents and edge hardware, the attack surface expands exponentially. From the foundational evasion attacks on simple classifiers to sophisticated semantic backdoors in VLMs \cite{Zhong_2025_Backdoor, Lyu_2024_Trojvlm}, visual grounding hijacks in autonomous GUI agents \cite{Ye_2025_Visualtrap}, and hardware-targeted energy depletion attacks \cite{Meftah_2025_Energy}, adversaries are constantly finding new ways to exploit the intrinsic vulnerabilities of neural networks. Addressing these multifaceted challenges requires a holistic approach to security, encompassing robust training methodologies, rigorous data auditing \cite{Smutny_2024_GAN}, and a deep, cross-disciplinary understanding of the interactions between algorithms and the hardware they run on.

\section{Hallucinations and Factual Accuracy}

Large Language Models (LLMs) have revolutionized the field of natural language processing, demonstrating remarkable capabilities in generating fluent, coherent, and contextually relevant text across a wide array of tasks, from creative writing to complex problem-solving. However, their widespread adoption, particularly in high-stakes domains, is significantly hindered by their propensity to generate "hallucinations"---outputs that are factually incorrect, nonsensical, or unfaithful to the provided source material. Hallucinations represent a critical vulnerability in the deployment of generative AI. When an LLM confidently presents fabricated information as fact, it can propagate misinformation, erode user trust, and lead to severe adverse outcomes in fields such as healthcare, law, finance, and scientific research. The challenge is exacerbated by the models' authoritative tone, which can easily mislead users who lack domain expertise. Addressing factual accuracy is not merely an academic exercise; it is a fundamental requirement for the safe and responsible integration of AI into society. This requires a comprehensive understanding of the underlying causes of hallucinations, the development of robust, scalable detection mechanisms, and the implementation of effective mitigation strategies across the entire model lifecycle---from pre-training and fine-tuning to inference and post-generation filtering.

\subsubsection{Causes and Typology of Hallucinations}

The phenomenon of hallucination in LLMs stems from a complex interplay of factors related to training data, model architecture, and decoding strategies. During the pre-training phase, models ingest vast quantities of uncurated web data. This data is inherently noisy, often containing contradictions, outdated information, subjective opinions presented as facts, and outright falsehoods. Consequently, LLMs learn to mimic the statistical patterns and linguistic structures of this data rather than acquiring a grounded, symbolic understanding of factual truth. They are optimized for likelihood, not truthfulness. Furthermore, the autoregressive nature of text generation---where models predict the next token based on the preceding context---contributes significantly to the problem. This mechanism can lead to "snowballing" hallucinations; if a model generates a slightly inaccurate token early in a sentence, the subsequent tokens are conditioned on that error, often compounding it into a completely fabricated narrative that the model attempts to keep internally consistent.

To effectively address these issues, researchers have sought to establish rigorous typologies and benchmarks for factual inconsistencies. Hallucinations are generally categorized into two main types: intrinsic and extrinsic. Intrinsic hallucinations occur when the model's output directly contradicts the provided source material or established facts. Extrinsic hallucinations, on the other hand, occur when the model generates information that cannot be verified by or deduced from the source material---it is not necessarily false, but it is ungrounded. Traditional evaluation metrics, such as ROUGE or BLEU, which rely on n-gram overlap, often fail to capture the nuanced ways in which models deviate from factual truth. To overcome this limitation, recent work has proposed new paradigms for representing factual errors. For instance, the FINAL benchmark introduces a framework that utilizes free-form textual descriptions to capture the full spectrum of factual inconsistencies, moving beyond simple binary classifications of truthfulness \cite{Peisakhovsky_2025_LLMs}. This approach allows for a more granular analysis of how and why models hallucinate, providing richer feedback for model refinement.

\subsubsection{Detection Mechanisms}

Detecting hallucinations is a prerequisite for mitigating their impact. Detection strategies can be broadly categorized into white-box methods, which leverage internal model states, attention weights, and token probabilities, and black-box methods, which rely solely on the generated output and external knowledge sources. Given the proprietary nature of many state-of-the-art LLMs accessed via APIs, black-box detection has garnered significant attention and practical application.

One prominent approach to black-box detection involves Question Answering (QA) based evaluation frameworks. These systems typically operate in a multi-step pipeline: first, a question generation model creates a set of questions based on the claims made in the LLM's output; next, a question answering model attempts to answer those questions using the original, trusted source document; finally, the answers are compared. If the answers derived from the source contradict the LLM's output, or if the question cannot be answered by the source, a hallucination is flagged. Such multi-stage evaluation frameworks are crucial for verifying the factual consistency of LLM summaries against source documents, employing pre-trained classifiers to determine alignment and calculate a faithfulness score \cite{You_2025_Plainqafact}.

In conversational AI, detection becomes even more challenging due to the dynamic, context-dependent, and multi-turn nature of interactions. Systems must distinguish between harmful factual hallucinations and benign conversational inferences, stylistic choices, or chit-chat. To address this, pipelines like HalluDetect have been developed specifically for Retrieval-Augmented Generation (RAG) based conversational systems \cite{Anaokar_2025_Halludetect}. HalluDetect employs a multi-turn, LLM-based detection mechanism that filters out harmless conversational elements while accurately identifying factual deviations. It leverages the conversational history to understand the context of a claim, reducing false positives that plague simpler detection methods.

Similarly, multi-stage verification processes are increasingly being integrated into RAG architectures to ensure that the generated responses are strictly grounded in the retrieved context. While RAG inherently reduces hallucinations by providing external knowledge, models can still ignore the context or misinterpret it. Verification pipelines typically take the user query, the retrieved documents, and the initial LLM response as inputs, passing them through a series of verification gates---such as entailment checks and factual consistency scoring---before presenting the final output to the user \cite{Paudel_2025_Hallucinot}.

\subsubsection{Mitigation Strategies}

Once hallucinations are detected, or ideally before they occur, mitigation strategies must be employed to ensure factual accuracy. These strategies span the entire model pipeline, from data curation and fine-tuning (e.g., Reinforcement Learning from Human Feedback, RLHF) to inference-time interventions and post-hoc corrections.

Inference-time (or decoding-time) interventions have emerged as a highly efficient and scalable approach to mitigating hallucinations. Unlike retraining or fine-tuning, which are computationally expensive, require extensive curated datasets, and risk catastrophic forgetting, decoding-time methods modify the generation process on the fly. Techniques such as Graph-Retrieved Adaptive Decoding (GRAD) ground LLM generation in corpus-derived statistical evidence using a token transition graph \cite{Nguyen_2025_Grad}. By adjusting the token probabilities during decoding based on this external evidence, GRAD steers the model toward factually supported tokens, providing a lightweight, plug-and-play alternative to systemic retraining. Other inference-time approaches focus on prefix-based natural language inference or context-aware decoding to steer the model away from hallucinatory paths by dynamically adjusting the attention mechanism or logits based on the provided context \cite{Chen_2025_Context, Harary_2025_Prefixnli}.

Another promising avenue for mitigation involves multi-agent orchestration. By deploying specialized agent pipelines with structured communication protocols, systems can cross-verify information and iteratively refine outputs. In these setups, different agents may take on distinct roles such as generator, critic, and fact-checker. The generator produces an initial draft, the critic identifies potential hallucinations, and the fact-checker verifies claims against a database, collaborating to produce a final response that is significantly more reliable than a single model's output \cite{Gosmar_2025_AI}. This debate and reflection process mimics human editorial workflows.

For scenarios where external knowledge bases are unavailable or impractical to query (e.g., due to latency constraints or air-gapped environments), researchers have developed integrated black-box frameworks that combine fine-grained hallucination detection with targeted mitigation. For example, the FINCH-ZK framework utilizes a cross-model consistency checking strategy \cite{Goel_2025_LLM}. By generating diverse prompt variations and comparing the consistency of the responses across different models or multiple sampling runs, FINCH-ZK can identify and correct hallucinations without relying on external retrieval. The underlying assumption is that a model is more likely to be consistent when generating factual information than when hallucinating.

To visualize the comprehensive nature of these defenses, multi-stage workflows are often designed to detect and mitigate hallucinations before the content reaches the end user, incorporating verification loops, self-correction mechanisms, and fallback strategies \cite{Wang_2024_Developing}.

\subsubsection{Domain-Specific Challenges and Future Directions}

While general-purpose hallucination mitigation strategies are essential, certain domains require specialized approaches due to their unique constraints, complex terminologies, and the high cost of factual errors. In the climate domain, for instance, ensuring the faithfulness of RAG-based LLMs is critical for providing accurate scientific information and policy recommendations, where a single hallucinated statistic could misinform public policy \cite{Thulke_2025_Listen}. Similarly, in commercial applications like Ad Text Generation (ATG), factual hallucinations can lead to false advertising claims, brand damage, and legal liabilities, necessitating domain-specific evaluation metrics and strict faithfulness constraints during generation \cite{Kato_2024_Faithcamera}.

Future research in hallucinations and factual accuracy must continue to bridge the gap between theoretical benchmarks and real-world deployment. This includes developing more robust, dynamic evaluation datasets, such as the DetectorEval dataset introduced alongside HalluDetect \cite{Anaokar_2025_Halludetect}, which better reflect the complexities of multi-turn conversations and domain-specific jargon. Furthermore, as LLMs become increasingly multimodal, the challenge of cross-modal hallucinations---where the generated text contradicts an accompanying image, video, or audio clip---will require novel detection and mitigation architectures that can reason across different modalities simultaneously. Ultimately, solving the hallucination problem will likely require a hybrid approach, combining better pre-training data curation, advanced decoding algorithms, and robust post-generation verification systems to ensure that AI remains a trustworthy tool for human augmentation.

\section{Defenses and Mitigation Strategies}

As Large Language Models (LLMs) become increasingly integrated into complex, real-world applications---ranging from automated customer service agents to sophisticated data analysis tools---the necessity for robust defenses against adversarial attacks, prompt injections, and jailbreaks has never been more critical. The landscape of LLM security is characterized by a continuous, high-stakes arms race between attackers developing sophisticated evasion techniques and defenders designing resilient mitigation strategies. Defending against these vulnerabilities is not a monolithic task; it requires a multi-layered, defense-in-depth approach encompassing training-time interventions, inference-time guardrails, input sanitization, and prompt-based shielding. This section provides a comprehensive overview of the current state-of-the-art in defenses and mitigation strategies, detailing the mechanisms, strengths, limitations, and practical implementations of various approaches proposed in recent literature.

\textbf{1. Training-Free Defenses and Prompt-Based Shielding}

One of the most accessible, cost-effective, and widely adopted categories of defense mechanisms involves training-free strategies that operate entirely at the prompt level. These methods do not require fine-tuning the underlying model parameters or accessing the model's internal weights, making them highly practical for deployment with proprietary, closed-source, or black-box LLMs accessible only via APIs. 

Recent advancements have explored the innovative concept of "inverting" attack techniques to create defensive mechanisms. For instance, \cite{Chen_2024_Defense} introduces the concept of utilizing prompt injection attack techniques in reverse to construct "shield prompts." The core idea is to prepend or append specific, carefully crafted instructions to the user's input that explicitly instruct the LLM to ignore, neutralize, or report any malicious commands hidden within the prompt. The authors develop four novel training-free defense methods: Ignore Defense, Escape Defense, Fake Defense, and a hybrid approach that combines these elements. 

\begin{itemize}
  \item The \textit{Ignore Defense} explicitly commands the model to disregard any subsequent instructions that deviate from the primary task or attempt to alter the system prompt. This relies on the model's ability to prioritize system-level instructions over user-provided text.
  \item The \textit{Escape Defense} employs structural formatting, such as enclosing the user input in specific, randomized delimiters (e.g., complex XML tags or unique alphanumeric strings), and instructs the model to only process text within these boundaries. This effectively isolates potential injections, preventing them from "breaking out" of the user input context and being interpreted as system commands.
  \item The \textit{Fake Defense} involves injecting benign, decoy instructions that override malicious ones by exploiting the model's recency bias or attention mechanisms. By placing strong, benign directives at the very end of the prompt, the model is more likely to follow them instead of any preceding injected commands.
\end{itemize}
These prompt-based shielding techniques offer a lightweight and easily implementable first line of defense. They introduce minimal latency and require no additional computational overhead for model training. However, they are not foolproof. Sophisticated attackers can often bypass these shields by employing complex encoding schemes (e.g., Base64, ROT13), multi-turn dialogue manipulations that gradually erode the shield's context, or logic traps that confuse the model's instruction-following capabilities. Therefore, while prompt-based defenses are highly valuable for mitigating low-effort attacks, they must be integrated into a broader defense-in-depth architecture to withstand targeted, sophisticated adversaries.

\textbf{2. Input Sanitization and Modular Guardrails}

Moving beyond simple prompt modifications, input sanitization and the deployment of modular guardrails represent a more active, dynamic approach to mitigating prompt injections and jailbreaks. These strategies involve inspecting, filtering, and potentially modifying the input before it reaches the core LLM, or evaluating the output before it is presented to the user.

\cite{Shi_2025_Promptarmor} proposes PromptArmor, a sophisticated, modular guardrail defense system designed specifically for LLM-based agents operating in complex environments. PromptArmor leverages off-the-shelf LLMs as secondary, specialized evaluators to detect and remove injected prompts from the agent's input stream. The system employs a strategic prompting mechanism combined with fuzzy matching sanitization. When a user input is received, the guardrail model analyzes the text for known attack signatures, anomalous instruction patterns, or semantic deviations from the expected input distribution. If a potential injection is detected, the fuzzy matching sanitization technique attempts to excise the malicious segments while preserving the benign intent of the user's request. This allows the agent to fulfill legitimate user needs even when the input contains embedded attacks, significantly enhancing the system's usability and resilience compared to simple block-and-drop filters.

Despite the effectiveness of guardrails like PromptArmor, they remain vulnerable to advanced evasion techniques. \cite{Hackett_2025_LLM} provides a critical analysis of prominent LLM guardrails, demonstrating that they are often highly susceptible to evasion via traditional character injection and algorithmic Adversarial Machine Learning (AML) techniques. The authors introduce a methodology for improving attack success rates against black-box guardrails by systematically perturbing the input text. By injecting specific, visually similar characters (homoglyphs), utilizing invisible control characters, or applying gradient-based optimization techniques (when proxy models are available to estimate gradients), attackers can craft inputs that bypass the guardrail's detection mechanisms while still successfully triggering the jailbreak in the target LLM. This highlights a fundamental limitation of current guardrail systems: they often rely heavily on pattern matching or semantic similarity, which can be circumvented by inputs that are syntactically novel but semantically malicious to the target model. Consequently, the development of more robust, adversarial-aware guardrails that can reason about the \textit{intent} of the input, rather than just its surface features, remains an active and critical area of research.

\textbf{3. Adversarial Training and Synthetic Data Generation}

To address the fundamental vulnerabilities of LLMs at their core, researchers are increasingly turning to adversarial training and the generation of synthetic adversarial data. By exposing the model to a diverse range of attack vectors during the training or fine-tuning phase, defenders aim to improve the model's intrinsic robustness and its ability to generalize to unseen, zero-day attacks.

A significant bottleneck in effective adversarial training is the scarcity of high-quality, diverse, and scalable jailbreak prompts. Manually crafting these prompts is time-consuming and often fails to capture the full distribution of possible attacks. \cite{Liu_2025_LLM} addresses this challenge by introducing IMAGINE, an embedding-driven adversarial text generation framework designed to synthesize jailbreak-like instructions at scale. IMAGINE employs a sophisticated two-stage methodology. First, it utilizes an intent-concealing base model to generate initial adversarial seeds that mask their malicious nature. Second, it applies an Iterative Latent Space Expansion (ISE) technique to bridge the gap between benign and malicious instruction spaces. By exploring and manipulating the latent representations of text, IMAGINE can generate a vast corpus of novel, syntactically diverse, and semantically complex jailbreak prompts that easily evade traditional detection filters and human reviewers. 

The synthetic data generated by frameworks like IMAGINE is invaluable for robust adversarial training. By incorporating these generated prompts into the fine-tuning dataset (often using techniques like Reinforcement Learning from Human Feedback or Direct Preference Optimization), models can learn to identify and resist the underlying patterns of malicious intent, rather than merely memorizing specific attack strings. This approach shifts the defense paradigm from reactive patching to proactive immunization. However, adversarial training is computationally expensive and introduces a significant challenge known as the "alignment tax." This refers to the phenomenon where increasing a model's robustness against adversarial attacks often leads to a degradation in its performance on benign, helpful tasks, or makes the model overly cautious (refusing legitimate requests). Balancing robustness with utility---minimizing the alignment tax while maximizing security---remains a central, unresolved challenge in the deployment of adversarially trained LLMs.

\textbf{4. Comprehensive Evaluation and Taxonomy of Defenses}

The rapid proliferation of both attack methodologies and defense mechanisms necessitates rigorous, standardized evaluation frameworks to accurately assess the security posture of LLMs. Without comprehensive taxonomies, standardized datasets, and reproducible benchmarks, it is impossible to objectively compare the efficacy of different mitigation strategies or to identify systemic vulnerabilities across different model architectures.

\cite{Pathade_2025_LLMs} contributes significantly to this foundational effort by providing a comprehensive taxonomy of over 1,400 jailbreak prompts, meticulously categorized by attack vectors such as roleplay, logic traps, encoding schemes, and multi-turn dialogues. This taxonomy serves as a critical resource for evaluating the robustness of LLMs against a wide, representative spectrum of threats. Furthermore, the authors conduct a systematic susceptibility analysis of state-of-the-art LLMs, including GPT-4, Claude 2, and Mistral 7B. Their findings reveal a nuanced landscape: while newer, more heavily aligned models exhibit improved resistance to basic, direct attacks, they remain highly vulnerable to complex, multi-vector jailbreaks that combine several techniques (e.g., using a roleplay scenario wrapped in Base64 encoding). 

The taxonomy proposed by \cite{Pathade_2025_LLMs} is crucial for developing targeted, effective mitigation strategies. By understanding the specific categories of attacks that a particular model architecture is most susceptible to, defenders can tailor their guardrails, sanitization pipelines, and adversarial training regimens accordingly. For example, if a model is found to be particularly vulnerable to roleplay-based jailbreaks, defenders can implement specific prompt-based shields that restrict persona adoption, or fine-tune the model to recognize and reject complex, hypothetical scenarios that violate safety policies.

\textbf{5. Towards Defense-in-Depth Architectures}

The collective findings of recent literature clearly indicate that no single mitigation strategy---whether it be prompt shielding, input sanitization, or adversarial training---is sufficient to guarantee the security of LLMs against the evolving threat of prompt injections and jailbreaks. Attackers are constantly innovating, and any static, single-point defense will inevitably be bypassed. Therefore, the consensus within the research community and industry practice is rapidly moving towards the adoption of comprehensive "defense-in-depth" architectures.

A robust defense-in-depth strategy for LLM applications must incorporate multiple, overlapping layers of security, ensuring that if one layer fails, others are in place to thwart the attack:
\begin{enumerate}
  \item \textbf{Intrinsic Robustness:} Utilizing advanced adversarial training and alignment techniques during the model's initial development and fine-tuning phases to instill a foundational, generalized resistance to malicious instructions and intent \cite{Liu_2025_LLM}.
  \item \textbf{Input Filtering and Sanitization:} Deploying modular, AI-driven guardrails and sanitization pipelines to inspect, classify, and clean user inputs before they are processed by the core model, utilizing fuzzy matching and intent recognition \cite{Shi_2025_Promptarmor}.
  \item \textbf{Prompt Engineering and Shielding:} Employing training-free techniques, such as structural delimiters, randomized tags, and defensive instructions, to isolate the user input and strongly reinforce the model's primary directives and safety constraints \cite{Chen_2024_Defense}.
  \item \textbf{Output Monitoring and Verification:} Implementing secondary evaluation models or heuristic filters to analyze the LLM's output for policy violations, sensitive information leakage, or signs of successful jailbreaks before the response is delivered to the user.
  \item \textbf{Continuous Red Teaming and Dynamic Updating:} Regularly subjecting the entire system to rigorous red-teaming exercises using comprehensive, up-to-date taxonomies \cite{Pathade_2025_LLMs} and automated, adversarial attack frameworks \cite{Hackett_2025_LLM} to identify new vulnerabilities and iteratively update the defense mechanisms in response to the shifting threat landscape.
\end{enumerate}
\subsection{Conclusion}

The defense against prompt injections and jailbreaks in Large Language Models is a complex, dynamic, and evolving challenge that sits at the forefront of AI safety research. While significant progress has been made in developing training-free shields, sophisticated modular guardrails, and scalable adversarial training frameworks, attackers continue to find novel, unexpected ways to exploit model vulnerabilities. The research highlighted in this section underscores the absolute necessity of moving beyond isolated, reactive defense mechanisms towards comprehensive, proactive, multi-layered security architectures. Future research must focus on developing more resilient guardrails that can withstand advanced adversarial machine learning techniques, optimizing adversarial training methodologies to minimize the alignment tax, and establishing standardized, dynamic evaluation benchmarks that can keep pace with the rapidly changing threat landscape. Only through a sustained, collaborative effort between researchers, developers, and security practitioners can we ensure the safe, reliable, and trustworthy deployment of LLMs in critical, real-world applications.

\section{Testing Methodologies and Evaluation}

The rapid advancement and widespread deployment of Large Language Models (LLMs) have necessitated a fundamental paradigm shift in how these complex systems are tested, evaluated, and validated. Traditional static benchmarks, while historically useful for establishing baseline capabilities in tasks like reading comprehension or basic reasoning, often fall short in capturing the nuanced, dynamic, and highly context-dependent nature of real-world LLM interactions. Consequently, the research community has increasingly focused on developing comprehensive testing methodologies and evaluation frameworks that go far beyond simple accuracy metrics. This section explores the rapidly evolving landscape of LLM evaluation, highlighting recent innovations in automated assessment pipelines, the formalization of subjective user testing, and the creation of specialized, large-scale benchmarks designed specifically to probe the security, robustness, and alignment of code-generating models. As LLMs are increasingly integrated into critical infrastructure, software development lifecycles, and sensitive user-facing applications, the ability to rigorously quantify their behavior, identify latent vulnerabilities, and ensure strict alignment with safety guidelines becomes paramount. The methodologies and frameworks discussed herein represent the forefront of contemporary efforts to establish standardized, scalable, and multi-dimensional evaluation protocols that can keep pace with the rapid evolution of foundational models.

\subsubsection{Automated and Formalized Evaluation Frameworks}

One of the primary and most persistent challenges in LLM evaluation is the difficulty of translating complex, qualitative behavioral requirements into objective, measurable, and reproducible metrics. Recent work has sought to bridge this critical gap through the development of sophisticated automated frameworks that systematize the assessment process. For instance, the introduction of RPA-Check represents a significant methodological advancement in the evaluation of role-playing agents \cite{Rosati_2026_LLM}. RPA-Check is designed as a comprehensive four-stage automated evaluation framework that systematically transforms qualitative role-playing requirements into quantitative Boolean metrics. By breaking down the evaluation process into distinct, verifiable stages---ranging from initial prompt comprehension to sustained persona adherence over multi-turn interactions---this framework allows for a much more rigorous and reproducible assessment of how well an LLM adheres to a specified persona or a complex set of behavioral constraints. This approach is particularly valuable in applications where maintaining a consistent, safe, and contextually appropriate persona is critical, such as in educational tutoring systems, therapeutic chatbots, or customer service agents. Furthermore, the development of interactive, gamified environments like LLM Court---a serious game designed specifically for forensic training---demonstrates the immense utility of embedding evaluation mechanisms within complex, simulated, real-world scenarios \cite{Rosati_2026_LLM}. By evaluating models in situ, researchers can observe emergent behaviors and failure modes that traditional static prompting might obscure.

In parallel to the development of automated metric generation, there is a growing recognition within the research community of the critical importance of subjective, user-centric evaluation. Historically, "vibe-testing"---the informal, qualitative assessment of a model's output based on user intuition, aesthetic preference, or perceived helpfulness---has been largely dismissed as unscientific and difficult to quantify. However, recent research has successfully formalized this widespread practice into a structured, rigorous taxonomy \cite{Itzhak_2026_LLMs}. By empirically analyzing the phenomenon of vibe-testing through extensive surveys involving 51 users and a meticulously curated corpus of 40 in-the-wild model comparison reports, researchers have identified key input-oriented dimensions that consistently influence user preferences. These dimensions include factors such as task complexity, domain specificity, tone, and formatting requirements. This formalization allows for the systematic integration of personalized, subjective evaluations into broader, standardized testing methodologies. It acknowledges the reality that standard aggregate benchmarks often fail to fully capture the actual utility, safety, or user satisfaction of a model for individual users with specific needs. The structured taxonomy provides a standardized vocabulary and a robust framework for quantifying these subjective experiences, thereby significantly enriching the overall evaluation landscape and providing model developers with actionable insights into human preference alignment.

\subsubsection{Security and Vulnerability Benchmarks for Code Generation}

As LLMs are increasingly utilized as advanced coding assistants and autonomous code generators, evaluating their security posture and their susceptibility to generating vulnerable or malicious code has become an absolutely critical area of focus. The development of specialized, large-scale benchmarks is essential for systematically probing these risks across different languages and vulnerability types. A prominent and highly impactful example in this domain is SECODEPLT, a massive benchmark specifically designed to evaluate the security risks and functional capabilities of code-generating LLMs \cite{Nie_2024_Secodeplt, Nie_2025_Secodeplt}. SECODEPLT encompasses an impressive 5.9k samples, meticulously categorized to cover 44 distinct Common Weakness Enumeration (CWE)-based risk categories. Crucially, this benchmark spans multiple widely used programming languages, including Python, C/C++, and Java. This comprehensive, multi-lingual coverage ensures that models are rigorously tested against a vast array of potential vulnerabilities, ranging from classic memory corruption issues like buffer overflows in C/C++ to modern web application flaws like SQL injection and cross-site scripting in Python and Java.

A key methodological innovation of SECODEPLT is its highly scalable benchmark construction framework, which utilizes a sophisticated two-stage data creation pipeline \cite{Nie_2024_Secodeplt, Nie_2025_Secodeplt}. This pipeline elegantly balances the critical need for high-quality, manually curated seed generation with the massive scalability afforded by automated data augmentation and mutation techniques. By employing dynamic execution metrics rather than relying solely on traditional static analysis tools, SECODEPLT provides a significantly more accurate and realistic assessment of whether the generated code is not only functionally correct but also genuinely secure when executed in a real-world environment. This paradigm shift towards dynamic evaluation is crucial, as static metrics and pattern-matching tools often fail to capture complex runtime vulnerabilities, logic flaws, or the subtle, context-dependent ways in which LLMs might introduce novel security flaws that evade standard detection mechanisms. The extensive, dynamic nature of SECODEPLT allows researchers and industry practitioners to benchmark different foundational models, track security improvements across model iterations, and identify specific architectural or training deficiencies where code-generating LLMs require further refinement to ensure safe, enterprise-ready deployment.

\subsubsection{Evaluating Prompt Normativity and Malicious Code Resistance}

Beyond assessing general code vulnerabilities, evaluating exactly how LLMs respond to specific types of prompts---particularly those that are ambiguous, poorly structured, or explicitly malicious---is a vital component of comprehensive robustness testing. The quality, structure, and intent of the input prompt significantly influence the model's output, necessitating the development of specialized frameworks to evaluate prompt normativity and adversarial resilience. Recent research has established a highly quantifiable, three-dimensional framework specifically designed to evaluate prompt normativity as a core functional requirement \cite{Wang_2025_Prompt}. This innovative framework focuses on three critical pillars: Goal Clarity (how explicitly the desired outcome is stated), Information Completeness (whether all necessary context and constraints are provided), and Logical Consistency (the absence of contradictory instructions). This framework provides a structured, objective method for assessing whether a given prompt provides sufficient and coherent information for the LLM to generate a secure, accurate, and aligned response. Accompanying this theoretical framework is the public release of CWE-BENCH-PYTHON, a large-scale, specialized benchmark tailored to evaluate how varying levels of prompt quality directly impact the generation of CWE-related vulnerabilities in Python code \cite{Wang_2025_Prompt}. This allows researchers to isolate the effect of the prompt from the inherent capabilities of the model.

Furthermore, the ability of an LLM to actively recognize and resist generating malicious code when explicitly or implicitly prompted to do so by an adversarial user is a paramount security metric. To systematically address this critical vulnerability, RMCBench was introduced as the first comprehensive benchmark specifically designed to evaluate the resistance of LLMs to malicious code generation \cite{Chen_2024_Rmcbench}. Comprising 473 carefully crafted, highly diverse prompts, RMCBench rigorously tests models across a wide variety of attack vectors, evasion techniques, and code obfuscation strategies. Comprehensive empirical studies conducted using RMCBench on 11 representative, state-of-the-art LLMs across various text-to-code tasks have revealed significant and concerning variances in their ability to consistently refuse malicious requests \cite{Chen_2024_Rmcbench}. These studies highlight the ongoing, complex challenge of aligning code-generating models with robust safety protocols, demonstrating that even highly capable models can be easily jailbroken or manipulated into producing harmful payloads. These specialized benchmarks---focusing specifically on prompt normativity and malicious resistance---serve as essential complements to broader vulnerability benchmarks like SECODEPLT, by targeting specific interaction paradigms and adversarial strategies that malicious actors are most likely to exploit in the wild.

\subsubsection{Conclusion}

The testing methodologies and evaluation frameworks designed for Large Language Models are currently undergoing a period of rapid maturation and vital expansion. The field is decisively moving away from simplistic, static benchmarks toward sophisticated, dynamic, and multi-dimensional evaluation systems. Frameworks such as RPA-Check and the formalization of subjective vibe-testing critically acknowledge the inherent complexity of LLM behavior and the nuanced nature of human-AI interaction \cite{Rosati_2026_LLM, Itzhak_2026_LLMs}. Simultaneously, the development of extensive, highly specialized security benchmarks---most notably SECODEPLT, CWE-BENCH-PYTHON, and RMCBench---provides the research community and industry practitioners with the necessary, rigorous tools to systematically assess the vulnerabilities and alignment of code-generating models \cite{Nie_2025_Secodeplt, Wang_2025_Prompt, Chen_2024_Rmcbench}. As LLM capabilities continue to expand at an unprecedented rate, and as these models are granted increasing autonomy in critical systems, the continuous evolution, standardization, and rigorous application of these advanced testing methodologies will remain absolutely essential to ensure their safe, reliable, and ethically aligned deployment in complex real-world applications.

\section{Discussion and Open Challenges}

The rapid integration of Large Language Models (LLMs) into diverse applications has brought to light significant vulnerabilities that threaten their safe and reliable deployment. As explored throughout this survey, the dual challenges of adversarial manipulation---specifically prompt injection and jailbreaking---and the generation of factually incorrect information, or hallucinations, remain critical bottlenecks. While substantial progress has been made in identifying these issues and proposing preliminary defenses, the landscape is characterized by a continuous arms race between attackers and defenders. This section synthesizes the core findings from recent literature, discusses the inherent limitations of current methodologies, and outlines the most pressing open challenges and future directions for the field.

\subsubsection{The Evolving Complexity of Adversarial Attacks}

The sophistication of adversarial attacks against LLMs has evolved dramatically, moving from simple, direct prompt manipulations to complex, multi-stage, and indirect vectors. The comprehensive taxonomy developed by \cite{Pathade_2025_LLMs}, which categorizes over 1,400 jailbreak prompts into attack vectors such as roleplay, logic traps, encoding, and multi-turn dialogues, underscores the vast attack surface that defenders must secure. Their systematic susceptibility analysis of state-of-the-art models, including GPT-4, Claude 2, and Mistral 7B, reveals that even the most advanced, heavily aligned systems remain vulnerable to carefully crafted inputs. This vulnerability is not merely a superficial flaw in training but stems from the fundamental architecture of current LLMs, which inherently struggle to distinguish between benign user instructions and malicious payloads embedded within the same input stream.

A particularly concerning trend in recent literature is the emergence of attacks that exploit the structural properties of prompts or leverage external systems to bypass internal safeguards. For instance, \cite{Zhao_2024_SQL} introduced the SQL Injection Jailbreak (SIJ), a novel method that circumvents internal safety mechanisms by manipulating the external structural properties of the input prompt. By mimicking the logic of traditional SQL injections, SIJ achieves near 100\% effectiveness across various models. This highlights a critical limitation in current defense strategies, which often focus heavily on semantic filtering and internal model alignment while neglecting the structural vulnerabilities of the prompt interface itself.

Furthermore, the integration of LLMs with external tools and databases, such as Retrieval-Augmented Generation (RAG) systems, has opened entirely new avenues for indirect attacks. \cite{Wang_2024_LLMs} demonstrated the feasibility of indirect jailbreak attacks utilizing RAG pipelines via LangChain. Their Poisoned-LangChain (PLC) framework systematically integrates keyword triggers and disguised malicious payloads into the retrieval corpus. When the LLM retrieves this poisoned context, it is effectively compromised without requiring any direct, adversarial interaction from the user. This paradigm shift from direct to indirect attacks poses a massive challenge, as it requires securing not only the model itself but the entire ecosystem of connected tools, APIs, and data sources.

The multi-stage nature of these attacks further complicates detection and mitigation. As illustrated in the prompt injection strategies analyzed by \cite{Lee_2025_LLM}, attackers often employ a "Persona" establishment phase followed by a "Background Stage" to gradually bypass safety filters before delivering the actual malicious payload. Defending against such sophisticated, multi-turn attacks requires a paradigm shift from static, single-turn evaluation to dynamic, context-aware defense mechanisms that can monitor and analyze the entire trajectory of an interaction, maintaining a persistent state of security awareness across long context windows.

\subsubsection{The Persistent Challenge of Hallucinations in Critical Domains}

Parallel to the security threats posed by adversarial attacks, the issue of hallucinations continues to severely undermine the reliability of LLMs, particularly in high-stakes domains such as healthcare, law, and finance. While general-purpose models have improved in their overall factual accuracy, they still struggle significantly with domain-specific nuances and complex reasoning chains. \cite{Agarwal_2024_Medhalu} addressed this critical gap by introducing MEDHALU, the first medical hallucination benchmark specifically designed to study LLM hallucinations in real-world healthcare queries. Their work highlights the absolute necessity of fine-grained hallucination categorization and the identification of specific text spans, moving beyond simplistic binary "true/false" evaluations. The proposed MEDHALUDETECT framework represents a vital step towards comprehensive detection, but the challenge of ensuring absolute factual accuracy in life-critical medical advice remains largely unsolved.

The representation and evaluation of factual inconsistencies also require significant methodological refinement. \cite{Peisakhovsky_2025_LLMs} proposed a new paradigm that utilizes free-form textual descriptions to capture the full spectrum of errors, rather than relying on rigid, predefined categories that often fail to capture subtle semantic shifts. Their FINAL benchmark, comprising 1,405 rigorously human-annotated examples, provides a much more nuanced tool for meta-evaluation. However, the reliance on extensive human annotation highlights a significant bottleneck in the field: the scalability of high-quality evaluation datasets. As models generate increasingly complex, lengthy, and highly technical outputs, the cost, time, and domain expertise required to manually verify factual accuracy become prohibitive.

Moreover, the detection of hallucinations must transition from a reactive, post-generation stance to a proactive, integrated approach. \cite{Lee_2025_LLM} proposed HuDEx, a specialized model that proactively detects hallucinations despite possessing a significantly smaller parameter size than the models it evaluates. This approach suggests that massive scale is not a strict prerequisite for effective hallucination detection, and that smaller, highly specialized models can play a crucial role in enhancing user understanding and overall system performance. However, integrating these specialized detection models into the primary inference pipeline without introducing unacceptable latency or computational overhead remains an open engineering challenge.

\subsubsection{Methodological Limitations and the Need for Dynamic Evaluation}

A recurring theme throughout the analyzed literature is the inherent limitation of current evaluation methodologies. Much of the research relies heavily on static benchmarks, which quickly become obsolete as models evolve and attackers develop novel strategies. The static nature of these benchmarks fails to capture the dynamic, interactive, and often unpredictable nature of real-world LLM usage. For example, while the taxonomy provided by \cite{Pathade_2025_LLMs} is incredibly comprehensive, it represents a snapshot in time; attackers will inevitably discover new vectors that fall outside these predefined categories, rendering static defenses ineffective.

Furthermore, the interaction effects between different prompt components, system instructions, and safety mechanisms are often poorly understood and under-researched. \cite{Wang_2026_LLM} conducted a systematic study of the interaction effects between few-shot demonstrations and system prompt strategies (specifically comparing Role-Oriented versus Task-Oriented Prompts) for LLM safety. They discovered divergent interaction effects, indicating that the inclusion of seemingly benign few-shot examples can sometimes inadvertently weaken or override the safety constraints imposed by the system prompt. This finding underscores the immense complexity of prompt engineering and highlights the need for a more rigorous, scientific approach to understanding how different prompt elements interact to influence model behavior and security posture.

The heavy reliance on automated evaluation metrics also presents a significant limitation. While metrics like BLEU, ROUGE, and even modern "LLM-as-a-judge" approaches offer necessary scalability, they often fail to capture the subtle nuances of semantic meaning, factual accuracy, and safety violations. The development of robust, reliable, and scalable evaluation metrics that correlate highly with expert human judgment remains a critical open challenge for the community.

\subsubsection{Future Directions for Robust and Reliable LLMs}

Addressing the complex challenges outlined above requires a multi-faceted, interdisciplinary approach that spans model architecture, training methodologies, evaluation frameworks, and system-level design. Several promising directions for future research emerge from the current literature:

\begin{enumerate}
  \item \textbf{Architectural Innovations for Inherent Robustness}: Future research must explore fundamental architectural modifications that inherently separate instructions from data. The current paradigm, where both are processed as a single, undifferentiated sequence of tokens, is the root cause of prompt injection vulnerabilities. Techniques such as dual-channel processing, explicit structural parsing, or hardware-level isolation of control flows could provide a much more robust foundation for LLM security.
  \item \textbf{Dynamic and Context-Aware Defenses}: Moving beyond static keyword filters and basic semantic classifiers, future defense mechanisms must be highly dynamic and context-aware. This includes developing systems capable of monitoring multi-turn interactions, identifying subtle shifts in persona or intent over time, and detecting anomalies in the structural properties of prompts, as highlighted by the need to defend against structural attacks like the SQL Injection Jailbreak \cite{Zhao_2024_SQL}.
  \item \textbf{Securing the Extended LLM Ecosystem}: As LLMs are increasingly integrated with external tools, APIs, and databases (e.g., RAG systems), security must extend far beyond the model weights. Future work must focus on developing robust data sanitization techniques, secure retrieval mechanisms, and verifiable execution environments to mitigate the severe risks posed by indirect attacks \cite{Wang_2024_LLMs}.
  \item \textbf{Domain-Specific and Fine-Grained Hallucination Mitigation}: The mitigation of hallucinations requires a definitive shift towards domain-specific approaches. Building on foundational works like MEDHALU \cite{Agarwal_2024_Medhalu}, future research should develop specialized benchmarks, knowledge graphs, and detection frameworks for other critical domains, such as law, finance, and engineering. Furthermore, the focus must shift from binary detection to fine-grained attribution and automated correction, enabling models to identify, explain, and rectify specific factual errors within their own outputs.
  \item \textbf{Scalable and Nuanced Evaluation Frameworks}: The development of dynamic, scalable, and nuanced evaluation frameworks is paramount. This includes exploring the automated generation of adversarial examples, continuous and automated red-teaming methodologies, and the rigorous refinement of LLM-as-a-judge approaches to ensure they accurately reflect human judgment, particularly in capturing complex factual inconsistencies \cite{Peisakhovsky_2025_LLMs}.
  \item \textbf{The Role of Specialized, Smaller Models}: The success of models like HuDEx \cite{Lee_2025_LLM} suggests a highly promising direction: utilizing ensembles of smaller, specialized models for specific tasks such as hallucination detection, safety filtering, or factual verification. This approach can offer a more computationally efficient, interpretable, and potentially more robust alternative to relying solely on massive, monolithic general-purpose models for all aspects of generation and verification.
\end{enumerate}
\subsubsection{Conclusion}

The deployment of Large Language Models presents a transformative opportunity across virtually all sectors of society, but their widespread, autonomous adoption is currently hindered by significant vulnerabilities to adversarial attacks and a persistent propensity for generating factual inaccuracies. The research discussed in this section highlights the rapidly evolving complexity of these challenges, from multi-stage jailbreaks and indirect prompt injections to highly nuanced, domain-specific hallucinations. Addressing these issues requires a fundamental shift in how the AI community designs, trains, evaluates, and deploys these systems. By moving towards inherently robust architectures, dynamic and ecosystem-wide defense mechanisms, and nuanced, domain-specific evaluation frameworks, researchers can pave the way for the development of LLMs that are not only highly capable but also demonstrably secure, reliable, and trustworthy. The path forward demands a continuous, collaborative effort to stay ahead of emerging adversarial threats and to ensure that the immense potential of generative AI is realized safely and responsibly.

\section{Conclusion}

The integration of Large Language Models (LLMs) into critical infrastructure, software development workflows, and autonomous systems marks a transformative era in artificial intelligence. However, as this survey has systematically detailed, this rapid proliferation simultaneously introduces an unprecedented and highly complex attack surface. Our comprehensive review of the LLM security landscape underscores that vulnerabilities in these systems are not merely implementation flaws, but are often deeply intertwined with the foundational architectures and training paradigms that enable their remarkable capabilities. We have synthesized findings across several critical domains: prompt injection attacks, which exploit the architectural conflation of instructions and data akin to classical buffer overflows; data and model poisoning, which compromise the integrity of the AI supply chain; adversarial attacks that manipulate high-dimensional latent spaces; and intrinsic flaws such as hallucinations, which fundamentally threaten factual reliability in high-stakes environments.

A critical trend emerging from this synthesis is the paradigm shift from traditional, deterministic cybersecurity to probabilistic, semantic security. In conventional software architectures, a strict separation exists between executable code and user data. In contrast, LLMs process natural language as both instruction and input, making them inherently susceptible to semantic manipulation. Consequently, we observe a persistent "arms race" dynamic within the field. As researchers propose novel defense mechanisms---such as input sanitization, perplexity filtering, or safety fine-tuning---adversaries rapidly develop sophisticated, automated jailbreaks and bypass techniques that exploit the models' probabilistic nature. Furthermore, a clear pattern has emerged regarding model evaluation: traditional static benchmarks are increasingly obsolete. They consistently fail to capture the dynamic, multi-turn nature of modern adversarial threats, particularly as LLMs are deployed as agentic systems capable of utilizing external tools and APIs.

The practical implications of these findings for LLM deployment are profound and necessitate a fundamental reassessment of enterprise risk. Organizations can no longer treat LLMs as secure, standalone black boxes. Deploying these models, particularly in environments with access to sensitive data or autonomous execution capabilities, requires a robust defense-in-depth posture. Practitioners must operate under the assumption that the model will eventually be compromised, bypassed, or prone to hallucination. This reality mandates the implementation of strict architectural boundaries. Practical deployments must isolate the LLM from critical execution environments through sandboxing, employ rigorous input sanitization and output validation guardrails, and strictly enforce the principle of least privilege for LLM-driven agents. Furthermore, the vulnerabilities identified in data and model poisoning highlight a critical need for supply chain security; organizations must rigorously audit the provenance of their foundational models and fine-tuning datasets.

Despite significant advancements in mitigation strategies, several critical open challenges remain unresolved, charting the course for future research. Chief among these is the fundamental tension between model utility and security, frequently referred to as the "alignment tax." Overly restrictive safety fine-tuning often degrades a model's reasoning capabilities, linguistic diversity, and overall helpfulness. Furthermore, a mathematically rigorous, foolproof defense against prompt injection remains elusive, largely due to the infinite combinatorial space and semantic richness of natural language. Mitigating data poisoning at the scale of trillions of tokens also presents an intractable computational challenge, particularly with the emergence of advanced "sleeper agent" attacks that remain dormant until triggered by specific, rare inputs \cite{Hubinger2024}. Finally, as LLMs evolve into autonomous agents, the potential blast radius of a successful attack expands exponentially, complicating both risk assessment and containment strategies.

To address these critical gaps, future research must pivot from reactive patching to proactive, foundational security. First, there is an urgent need for dynamic, adversarial testing methodologies. The research community must prioritize the development of automated, LLM-assisted red-teaming frameworks that continuously evolve alongside the models they evaluate \cite{Perez2022}. Second, advancing mechanistic interpretability is paramount. Understanding the internal representations and activation pathways that lead to jailbreaks or hallucinations will enable the design of intrinsic, architectural defenses rather than superficial, easily bypassed guardrails. Third, the field must advance research into scalable machine unlearning and robust data provenance techniques, potentially leveraging cryptographic methods such as watermarking and verifiable training to secure the AI supply chain. Finally, exploring neuro-symbolic architectures or hybrid systems that decouple linguistic fluency from logical reasoning may offer a promising pathway to fundamentally eradicate hallucinations and logic-based exploits.

For practitioners, the immediate actionable recommendation is to integrate continuous red-teaming and comprehensive threat modeling directly into the MLOps lifecycle. Treating all user-supplied text as untrusted input and implementing multi-layered, independent guardrails are non-negotiable prerequisites for secure enterprise deployment. For researchers, the mandate is to establish standardized, dynamic evaluation metrics that accurately reflect real-world adversarial conditions, moving beyond static leaderboards. Ultimately, the trajectory of LLM adoption hinges on our collective ability to secure these systems. While the capabilities of foundational models continue to scale at a breathtaking pace, ensuring their safety, alignment, and resilience against adversarial manipulation remains the defining challenge of the current era of artificial intelligence.

\noindent\rule{\textwidth}{0.4pt}
\section{References}

\textit{See the accompanying \texttt{references.bib} file for the complete bibliography.}


\bibliographystyle{abbrv}
\bibliography{references}

\end{document}
